\cleardoublepage
\thispagestyle{empty}
\begin{center}
	
	{\Large\textbf{Photon – Theory and Applications}}\\[1.2em]
	{\large Dipl.-Ing.\,(FH)\,Christian Weilharter}\\[1.2em]
	\textcopyright~2025, Christian Weilharter, Traunstein\\[2em]
\end{center}

\begin{flushleft}
	\begin{tabular}{@{}l l}
		\textbf{ISBN (Print):} & 978-3-912302-02-8 \\[0.5em]
		\textbf{ISBN (E-Book):} & 978-3-912302-03-5 \\[0.5em]
		\textbf{Edition:} & 1st edition, 2025 \\[0.5em]
		\textbf{Typesetting:} & \LaTeX \\[0.5em]
		\textbf{Publisher:} & Christian Weilharter \\[0.5em]

		\textbf{Print:} & Amazon KDP (Print-on-Demand) \\[0.5em]
		\textbf{E-book edition:} & Apple Books \\[0.5em]
		& Amazon Kindle Direct Publishing \\[0.5em]
		\textbf{Contact:} & \href{mailto:info@mathandphysics.de}{info@mathandphysics.eu} \\[0.5em]
		\textbf{Web:} & \href{https://www.mathandphysics.de}{www.mathandphysics.eu} \\

	\end{tabular}
\end{flushleft}

\vspace{2em}
\noindent
All rights reserved. No part of this book may be reproduced, stored, or transmitted
in any form or by any means, electronic or mechanical, including photocopying,
recording, or by any information storage and retrieval system, without prior written
permission of the author.

\begin{center}\small Printed in Germany\end{center}

\cleardoublepage


\chapter*{Preface}
\markboth{Preface}{Preface} % optional für Kopfzeile
% kein addcontentsline → Preface erscheint NICHT im Inhaltsverzeichnis



The creation of this book has been driven by a deep fascination with light and its fundamental mediator – the photon. In modern physics, the photon plays a central role: as a particle without mass, yet carrying energy and momentum; as the messenger of the electromagnetic interaction; and as a key figure in quantum mechanics.  
\noindent
My goal has been to present this multifaceted concept in its historical development, its physical significance, and its technical applications in such a way that both interested laypersons and advanced readers gain a solid understanding – without unnecessary oversimplification, yet always in a clear and accessible manner.  
\noindent
In the conception and preparation of this book, modern technology was also employed. The AI system ChatGPT by OpenAI (as of 2025) was used in a supportive role for formulating, structuring, and reflecting on the content. This collaboration made it possible to articulate ideas more quickly, test alternative phrasings, and clarify complex relationships.  
All content, however, was critically reviewed, revised, and is solely the responsibility of me.  
\noindent
Thus, this book is not only a contribution to the didactics of modern physics – it is also an experiment in how traditional science can gain a new level of clarity and accessibility through modern tools.  
\noindent
I hope that the enthusiasm for this subject will resonate with the readers – just as it has accompanied me for many years.  

\begin{flushright}
	\textit{Dipl.-Ing.(FH) Christian Weilharter} \\
	\vspace{0.5em}
	[Traunstein, 2025]
\end{flushright}
\newpage
\noindent
\chapter*{How the Book is Organized}
\markboth{How the Book is Organized}{How the Book is Organized} % no addcontentsline → this chapter will NOT appear in the table of contents

This book is divided into eight thematically coordinated chapters. It takes the reader from the physical foundations of light, through the quantization of the electromagnetic field, to modern applications and open questions in research. Each chapter highlights the role of the photon from a particular perspective—historical, theoretical, experimental, or technological.

\begin{itemize}
		\item \textbf{Chapter I} provides a brief overview of the historical development and theory of the photon.
	\item \textbf{Chapter II} describes the path from classical light theory to the introduction of the light quantum, covering central experiments such as the photoelectric effect\index{photoelectric effect} and blackbody radiation\index{blackbody radiation} that led to quantum theory.
	
	\item \textbf{Chapter III} addresses the physical properties of the photon, including energy, momentum, spin, polarization\index{polarization}, and masslessness.
	
	\item \textbf{Chapter IV} explores the experimental confirmation of the photon—from Millikan’s measurements\index{Millikan, Robert A.} of the photoelectric effect to modern single-photon experiments and quantum erasers\index{quantum eraser}.
	
	\item \textbf{Chapter V} introduces quantum electrodynamics (QED) and shows how the photon is understood as the gauge boson of the electromagnetic interaction.
	
	\item \textbf{Chapter VI} describes practical applications of the photon, such as in laser technology, medical imaging\index{medical imaging}, and communication technology.
	
	\item \textbf{Chapter VII} highlights current developments and future potential, particularly in quantum information\index{quantum information}, photonics, and fundamental research.
	
	\item \textbf{Chapter VIII} places the photon in the Standard Model of particle physics\index{Standard Model of particle physics}, showing how its masslessness arises from the symmetries of the theory and what open questions remain.
\end{itemize}
Numerous illustrations, quotations, and reflections accompany the scientific content, also shedding light on philosophical, historical, and epistemological dimensions. Mathematical tools are introduced carefully, so that the work remains accessible to interested readers with a general scientific background.

\begin{quote}
	\textit{With an understanding of the historical, conceptual, and experimental foundations of the photon, we now step into the terrain of quantization—the transition that carries light from the classical world into quantum physics.}
\end{quote}

\subsubsection*{How to Read This Book}
To clearly distinguish different aspects of the presentation, colored boxes are used throughout the book. They provide orientation and highlight central content:

\begin{itemize}
	\item \textbf{Physics boxes} (blue) provide physical explanations and background information.
	\item \textbf{Math boxes} (green) contain derivations, formulas, and mathematical details.
	\item \textbf{Didactics boxes} (yellow) point out conceptual pitfalls or offer alternative explanations.
	\item \textbf{Note boxes} (gray) provide structural hints or references to other chapters and appendices.
	\item \textbf{Warning boxes} (red) highlight particularly critical aspects or misunderstandings.
	\item \textbf{Hypothesis boxes} (orange) discuss hypothetical scenarios or “what-if” questions.
\end{itemize}

% -------------------------------------------------------------
\section*{Note on the Open-Access Edition}
% -------------------------------------------------------------

This book is part of the Open Science initiative \emph{“Christian \& Co-Pilot – Math \& Physics”}.

The complete, full-color PDF version is freely available at:

\begin{itemize}
	\item \href{https://mathandphysics.eu}{\texttt{https://mathandphysics.eu}}
	\item \href{https://zenodo.org/communities/christian-copilot}{\texttt{https://zenodo.org/communities/christian-copilot}}
\end{itemize}
The printed edition has been created for comfortable reading, long-term preservation,  
and to contribute to building an open-access library in print form.  
By purchasing this book, you support free scientific publishing  
and the idea of open and verifiable science.
